% ============================ CHAPITRE 1 =============================
\chapter{État de l'art des systèmes RAG multimodaux et positionnement du problème}
\label{chap:etat-art}

\section{Introduction}

Ce premier chapitre pose le cadre scientifique et industriel du projet. Il présente
d'abord le contexte de l'entreprise d'accueil et le domaine d'application (la recherche
d'information dans la documentation technique d'ingénierie), puis introduit les concepts
théoriques mobilisés : grands modèles de langue, représentations vectorielles
(\emph{embeddings}), génération augmentée par récupération (RAG) et modèles de
vision-langage (VLM). Il analyse ensuite les approches existantes pour rendre un corpus
multimodal interrogeable, les outils disponibles sur le marché et dans la littérature,
leurs limites, et conclut par le positionnement précis de la solution développée.

\section{Contexte : la documentation technique chez ALTEN}

\subsection{L'entreprise d'accueil}

ALTEN est un groupe international d'ingénierie et de conseil en technologies (ETC), présent
dans plus de trente pays. Le groupe accompagne les directions techniques et les directions
des systèmes d'information de grands acteurs industriels (aéronautique, automobile,
ferroviaire, énergie, télécommunications, sciences de la vie) sur l'ensemble du cycle de
développement de leurs produits et systèmes. Les activités du groupe couvrent notamment
l'ingénierie produit, les systèmes embarqués, les systèmes d'information industriels et,
de plus en plus, la valorisation des données et l'intelligence artificielle.

\subsection{Problématique métier}

Par la nature même de ses activités, ALTEN manipule quotidiennement une documentation
technique volumineuse et hétérogène : spécifications, manuels d'équipements, normes,
articles de recherche, guides méthodologiques. Cette documentation présente trois
caractéristiques qui rendent son exploitation difficile :

\begin{itemize}
  \item \textbf{Hétérogénéité des formats} : PDF issus de chaînes d'édition variées,
        documents Word, classeurs Excel, avec des mises en page complexes (colonnes
        multiples, tableaux imbriqués, en-têtes et pieds de page parasites) ;
  \item \textbf{Information visuelle critique} : dans un manuel d'ingénierie électrique,
        par exemple, la réponse à « que représente ce symbole ? » se trouve dans un
        \emph{schéma}, pas dans un paragraphe ; les diagrammes d'architecture des articles
        scientifiques concentrent souvent l'essentiel de la contribution ;
  \item \textbf{Coût d'accès élevé} : retrouver une information précise suppose de savoir
        dans quel document chercher, puis de le parcourir manuellement, un processus lent,
        frustrant et non scalable.
\end{itemize}

Le stage répond à cette problématique par une question d'ingénierie précise :
\emph{comment permettre à un ingénieur d'interroger en langage naturel un corpus de
documents techniques, textes et figures compris, et d'obtenir une réponse fiable,
vérifiable et illustrée ?}

\section{Concepts et notions théoriques clés}
\label{sec:concepts}

\subsection{Grands modèles de langue (LLM)}

Les grands modèles de langue sont des réseaux de neurones fondés sur l'architecture
\emph{Transformer} \cite{vaswani2017attention}, pré-entraînés sur d'immenses corpus à
prédire le token suivant. Cette architecture, introduite par Vaswani et al. en 2017,
remplace la récurrence par un mécanisme d'\emph{auto-attention} qui met en relation tous
les tokens d'une séquence en parallèle. Les modèles qui en découlent, de BERT
\cite{devlin2019bert}, encodeur bidirectionnel pré-entraîné par MLM, aux modèles génératifs
actuels (GPT-4o, Mistral, Llama, Gemma), présentent des capacités remarquables de
compréhension et de génération de texte. Utilisés seuls, ils souffrent cependant de deux
limites structurelles pour un usage documentaire : leur connaissance est figée à la date
d'entraînement et n'inclut pas les documents internes d'une organisation, et ils peuvent
produire des affirmations plausibles mais fausses (\emph{hallucinations}).

\subsection{Embeddings et recherche sémantique}

Un \emph{embedding} est une représentation d'un texte sous forme de vecteur dense de
dimension fixe (1024 dimensions pour \texttt{mistral-embed}, utilisé dans ce projet),
apprise de sorte que deux textes sémantiquement proches aient des vecteurs proches. La
similarité entre une question et un passage se mesure alors par similarité cosinus (ou par
distance euclidienne sur des vecteurs normalisés L2, ce qui revient au même classement).
La \emph{recherche sémantique} consiste à vectoriser le corpus une fois pour toutes, à
stocker ces vecteurs dans une \emph{base vectorielle}, puis, à chaque question, à
rechercher les $k$ vecteurs les plus proches du vecteur de la question. Contrairement à la
recherche lexicale, elle retrouve des passages pertinents même sans mot commun avec la
question (synonymes, reformulations, traductions).

\subsection{Recherche lexicale : BM25}

BM25 \cite{robertson2009bm25} est la fonction de classement lexical de référence,
héritière de TF-IDF : elle score un document selon la fréquence des termes de la requête
qu'il contient, pondérée par leur rareté dans le corpus et par la longueur du document.
Ses forces sont exactement complémentaires de celles des embeddings : elle excelle sur les
correspondances \emph{exactes} (identifiants, références normatives, « Figure 1 »,
acronymes), là où deux libellés proches (« Figure 1 » et « Figure 2 ») produisent des
embeddings quasi identiques que la similarité cosinus ne sait pas départager.

\subsection{Retrieval-Augmented Generation (RAG)}

Le RAG \cite{lewis2020rag} articule ces briques en un pipeline en deux temps
(figure~\ref{fig:rag-principe}) :

\begin{enumerate}
  \item \textbf{Récupération (retrieval)} : la question de l'utilisateur est vectorisée,
        puis les passages du corpus les plus pertinents sont récupérés depuis l'index ;
  \item \textbf{Génération augmentée} : ces passages sont injectés dans le prompt du LLM,
        avec l'instruction explicite de ne répondre \emph{qu'à partir d'eux} et de citer
        ses sources.
\end{enumerate}

Le LLM n'est donc plus une base de connaissances mais un \emph{moteur de synthèse
contraint} : la connaissance vit dans l'index documentaire, actualisable à tout moment,
et chaque affirmation de la réponse est traçable vers sa source.

\rapportfig[0.9]{rag_principe.png}{Principe général du paradigme RAG : indexation hors
ligne du corpus, puis récupération et génération contrainte à chaque question.}{fig:rag-principe}

\subsection{Le découpage sémantique (chunking)}

Les documents étant trop longs pour être vectorisés d'un bloc (et les fenêtres de contexte
des modèles d'embeddings limitées, typiquement 512 tokens), le corpus doit être découpé en
\emph{chunks}. La qualité de ce découpage conditionne directement celle du retrieval : un
découpage naïf à taille fixe coupe les phrases et mélange les sections, tandis qu'un
\emph{découpage sémantique} respecte la structure logique du document (sections,
paragraphes, tableaux) et rattache chaque chunk à son contexte hiérarchique (titres de
sections englobantes). C'est l'approche du \emph{HybridChunker} de Docling, utilisée ici.

\subsection{Modèles de vision-langage (VLM)}

Les modèles de vision-langage (GPT-4o vision, LLaVA, BLIP, Gemma-3) acceptent en entrée
une image et un texte, et produisent du texte. Appliqués à une figure extraite d'un
document, ils peuvent en produire une \emph{description textuelle détaillée} : type de
figure, composants, relations, textes et valeurs visibles. Cette capacité est la clé de
voûte du RAG \emph{multimodal} par description : elle projette le contenu visuel dans
l'espace du texte, où il devient indexable et interrogeable avec les mêmes outils que le
reste du corpus.

\section{Approches existantes pour le RAG multimodal}
\label{sec:approches}

La littérature et la pratique industrielle distinguent trois grandes stratégies pour
rendre les images d'un corpus interrogeables (figure~\ref{fig:approches-multimodal}) :

\subsection{OCR seul}
L'approche historique applique une reconnaissance optique de caractères (Tesseract, Azure
Document Intelligence) aux images pour en extraire le texte visible. Elle est efficace pour
les documents scannés et les captures d'écran riches en texte, mais ne capture rien de la
\emph{sémantique visuelle} : un schéma électrique passé à l'OCR ne livre que quelques
étiquettes isolées, pas la signification du circuit.

\subsection{Captioning par VLM (description textuelle)}
Chaque image est décrite par un VLM, et c'est la \emph{description} qui est vectorisée et
indexée aux côtés des chunks de texte, dans le même espace vectoriel. Cette approche offre
un excellent rapport qualité/complexité : index unique et homogène, pipeline de retrieval
inchangé, descriptions lisibles par l'utilisateur (utile pour l'affichage des sources).
Sa limite : la qualité de l'index dépend de la qualité (et de l'orientation) des
descriptions générées.

\subsection{Embeddings multimodaux natifs}
Des modèles comme CLIP \cite{radford2021clip} ou, plus récemment, ColPali
\cite{faysse2024colpali} projettent images et textes dans un espace vectoriel \emph{commun},
permettant de rechercher directement une image à partir d'une question textuelle, sans
description intermédiaire. C'est l'approche la plus puissante sur le plan théorique, mais
elle impose une infrastructure de retrieval spécifique (index séparé ou multi-vecteurs,
\emph{late interaction} pour ColPali), des modèles plus lourds à héberger, et produit des
représentations opaques pour l'utilisateur.

\rapportfig[0.95]{rag_multimodal_approches.png}{Les trois stratégies d'indexation des
images d'un corpus : OCR seul, description par VLM, embeddings multimodaux natifs.}{fig:approches-multimodal}

\section{Outils et solutions existants}
\label{sec:outils}

\subsection{Cadres logiciels open source}
\textbf{LangChain} et \textbf{LlamaIndex} sont les deux cadres dominants pour assembler
des pipelines RAG : ils fournissent des abstractions prêtes à l'emploi (loaders, splitters,
retrievers, chains). Leur généricité a un coût : couches d'abstraction épaisses, dépendances
nombreuses, comportements parfois difficiles à auditer, un point sensible lorsqu'il faut
diagnostiquer finement la qualité du retrieval.

\subsection{Analyse documentaire}
Pour l'extraction du contenu des PDF, l'écosystème s'étend des bibliothèques bas niveau
(PyMuPDF, pdfplumber, pypdfium2) aux solutions d'analyse structurelle par apprentissage
profond : \textbf{Docling} \cite{docling2024} (IBM Research, open source), qui combine un
modèle de détection de mise en page et TableFormer pour la reconstruction des tableaux, et
les services cloud managés comme \textbf{Azure Document Intelligence} ou Amazon Textract.

\subsection{Bases vectorielles}
Le marché offre un continuum allant des bibliothèques embarquées (\textbf{ChromaDB},
FAISS, LanceDB) aux serveurs dédiés (\textbf{Qdrant}, Weaviate, Milvus) et aux services
managés (Pinecone). Les premières s'exécutent dans le processus applicatif sans
infrastructure supplémentaire ; les seconds apportent réplication, sharding et très gros
volumes.

\subsection{Solutions intégrées du marché}
Des produits clés en main adressent le même besoin : NotebookLM (Google), les assistants
documentaires d'Azure AI Search « chat with your data », Glean, ou les GPTs personnalisés
d'OpenAI. Ils offrent une mise en œuvre rapide mais imposent l'envoi du corpus chez un
tiers, une personnalisation limitée du pipeline (chunking, retrieval, prompts) et des coûts
de licence : trois points bloquants pour un démonstrateur interne maîtrisé de bout en bout.

\section{Limites des solutions existantes}
\label{sec:limites-existant}

L'analyse de l'existant fait ressortir quatre limites au regard du besoin exprimé :

\begin{enumerate}
  \item \textbf{Cécité visuelle des RAG standards} : les pipelines RAG « par défaut »
        (et la plupart des tutoriels/outils grand public) n'indexent que le texte ; les
        figures, porteuses d'une part essentielle de l'information technique, sont ignorées ;
  \item \textbf{Retrieval purement dense} : la recherche par embeddings seuls échoue sur
        les requêtes exactes (références de figures, acronymes, identifiants), fréquentes
        dans un contexte d'ingénierie ;
  \item \textbf{Extraction PDF simpliste} : les loaders textuels basiques détruisent la
        structure (tableaux aplatis, colonnes mélangées, légendes orphelines), dégradant
        mécaniquement la qualité des chunks puis des réponses ;
  \item \textbf{Inadéquation aux environnements contraints} : les solutions supposant
        l'installation de services (Docker, serveurs de bases de données, OCR système) ou
        l'accès libre à Internet sont inapplicables sur un poste d'entreprise verrouillé
        derrière un proxy d'inspection TLS.
\end{enumerate}

\section{Apport de la solution proposée}
\label{sec:apport}

La solution développée répond point par point à ces limites :

\begin{keybox}[Apports de la solution]
\begin{itemize}
  \item \textbf{Indexation multimodale unifiée} : les figures sont extraites par Docling,
        décrites par un VLM (Gemma-3 27B) avec un prompt d'ancrage contextuel, et leurs
        descriptions sont indexées \emph{dans le même espace vectoriel} que le texte : une
        question peut donc être résolue indifféremment par un paragraphe ou par une figure,
        laquelle est alors affichée à l'utilisateur ;
  \item \textbf{Recherche hybride dense + BM25 fusionnée par RRF} : la complémentarité
        sémantique/lexicale est exploitée sans calibration de scores, par fusion de rangs
        réciproques \cite{cormack2009rrf} ;
  \item \textbf{Analyse structurelle de qualité} : Docling préserve la structure logique
        des documents (sections, tableaux, légendes de figures), et le découpage
        \emph{HybridChunker} produit des chunks contextualisés respectant la limite de
        tokens du modèle d'embeddings ;
  \item \textbf{Architecture sans dépendance système} : SQLite, ChromaDB persistant et
        FakeRedis s'exécutent dans le processus Python ; les modèles lourds sont consommés
        via des API cloud accessibles à travers le proxy d'entreprise ; le tout
        fonctionne sur un poste verrouillé et se déploie tel quel sur une VM OCI ;
  \item \textbf{Qualité mesurée, pas supposée} : le système est évalué quantitativement
        avec RAGAS \cite{es2023ragas} et un juge LLM indépendant du système évalué.
\end{itemize}
\end{keybox}

\section{Positionnement du projet}
\label{sec:positionnement}

Le tableau~\ref{tab:positionnement} synthétise le positionnement de la solution par rapport
aux familles d'alternatives. Le projet assume des choix de \emph{sobriété architecturale}
adaptés à son cadre (un démonstrateur de huit semaines destiné à un déploiement OCI léger) :
pipeline RAG assemblé directement sur les SDK des fournisseurs (sans LangChain), stockage
embarqué, approche multimodale par description VLM plutôt que par embeddings multimodaux
natifs. Conformément au cahier des charges, les architectures multi-agents, le RAG
auto-correctif et la gestion fine des droits sont explicitement hors périmètre et traités
comme perspectives.

\begin{table}[!htb]
\centering\small
\caption{Positionnement de la solution développée par rapport à l'existant.}
\label{tab:positionnement}
\begin{tabularx}{\textwidth}{@{}lXX@{}}
\toprule
\textbf{Dimension} & \textbf{Alternatives courantes} & \textbf{Choix du projet} \\
\midrule
Analyse PDF & Loaders textuels naïfs ; OCR cloud payant & Docling (layout + TableFormer), modèles locaux pré-téléchargés \\
Images & Ignorées, ou OCR seul & Captioning VLM contextualisé + indexation unifiée \\
Retrieval & Dense seul & Hybride dense + BM25, fusion RRF \\
Base vectorielle & Serveur dédié (Qdrant, Milvus) & ChromaDB embarqué persistant \\
Orchestration & LangChain / LlamaIndex & Pipeline explicite sur SDK natifs (auditable) \\
Génération & GPT-4o (OpenAI) & Mistral \texttt{mistral-small-latest} (accessible via le proxy, coût maîtrisé) \\
Évaluation & Tests manuels & RAGAS + juge LLM indépendant (Llama 3.3 70B) \\
\bottomrule
\end{tabularx}
\end{table}

\section{Conclusion}

Ce chapitre a établi le socle conceptuel du projet (LLM, embeddings, RAG, VLM, recherche
hybride) et montré que les solutions existantes laissent un espace précis : celui d'un
assistant documentaire multimodal, maîtrisé de bout en bout, fonctionnant dans un
environnement d'entreprise contraint. Le chapitre suivant traduit ce positionnement en une
architecture logicielle concrète : besoins, pipeline, composants, modélisation et
justification détaillée des choix technologiques.
